\section{Tổng kết}

\indent \indent Trong phần này, nhóm tổng kết lại những kết quả đã đạt được trong quá trình
thực hiện đồ án tốt nghiệp, đánh giá hiệu quả của hệ thống đã được hiện thực
hóa, cũng như nhận diện những hạn chế còn tồn tại và đề xuất hướng phát triển
tương lai cho hệ thống đặt vé và gợi ý sự kiện thông minh tích hợp chatbot AI.

% -----------------------------------------------------------
\subsection{Đánh giá kết quả hiện thực hệ thống giai đoạn Đồ án tốt nghiệp}

\indent \indent Sau quá trình nghiên cứu, thiết kế và triển khai, nhóm đã hoàn thiện một hệ
thống đặt vé sự kiện trực tuyến tích hợp mô hình gợi ý sự kiện thông minh và hệ thống trợ lý ảo - Chatbot AI, vận hành ổn định trên môi trường cloud với backend được
deploy trên \textbf{Render} và frontend được deploy trên \textbf{Vercel}.
Dưới đây là tổng hợp các kết quả đã đạt được theo từng nhóm chức năng.

\subsubsection*{(1) Quản lý tài khoản và phân quyền}

\indent \indent Hệ thống hiện thực thành công cơ chế xác thực dựa trên \textbf{JWT} với ba
vai trò hoàn toàn tách biệt: \textbf{user} (người dùng cuối), \textbf{organizer}
(ban tổ chức) và \textbf{admin} (quản trị viên). Mỗi vai trò được định tuyến
đến giao diện riêng biệt (\texttt{ClientLayout}, \texttt{OrganizerLayout},
\texttt{AdminLayout}) với bộ tính năng độc lập, đảm bảo nguyên tắc phân quyền
tối thiểu (\textit{principle of least privilege}). Luồng quên mật khẩu được
hiện thực hoàn chỉnh thông qua email tự động sinh mật khẩu mới và gửi đến
người dùng qua Nodemailer.

\subsubsection*{(2) Quản lý sự kiện và vé}

\indent \indent Ban tổ chức có thể tạo sự kiện với đầy đủ thông tin: tiêu đề, thể loại
(11 thể loại chuẩn hóa), mô tả dạng rich-text (Quill editor), poster, thời
gian, địa điểm (GeoJSON), sức chứa, giới hạn tuổi và sơ đồ chỗ ngồi. Mỗi
sự kiện hỗ trợ \textbf{nhiều hạng vé} (\texttt{TicketClass}) với hai loại
chỗ ngồi: tự do (\texttt{general}) và có ghế cụ thể (\texttt{reserved}),
kèm cơ chế kiểm soát số lượng và trạng thái bán vé (\texttt{available},
\texttt{sold\_out}, \texttt{unavailable}). Quản trị viên có toàn quyền duyệt,
chỉnh sửa và theo dõi toàn bộ sự kiện trên hệ thống.

\subsubsection*{(3) Luồng đặt vé và thanh toán}

\indent \indent Nhóm đã tích hợp thành công cổng thanh toán \textbf{PayOS} -- một cổng thanh
toán nội địa phổ biến tại Việt Nam hỗ trợ chuyển khoản ngân hàng và QR
thanh toán. Luồng đặt vé bao gồm các bước: chọn vé $\to$ nhập voucher
(tùy chọn) $\to$ tạo đơn với trạng thái \texttt{pending} $\to$ thanh toán
qua PayOS $\to$ xác nhận qua webhook $\to$ cập nhật trạng thái
\texttt{paid} $\to$ gửi email xác nhận kèm QR vé. Đơn hàng chưa thanh
toán sau 15 phút sẽ tự động bị thu hồi bởi cron job, đảm bảo slot vé
không bị chiếm giữ vô thời hạn. Hệ thống voucher hỗ trợ chiết khấu
theo số tiền cố định với điều kiện chi tiêu tối thiểu và ngày hết hạn.

\subsubsection*{(4) Vé điện tử và QR Code}

\indent \indent Mỗi vé sau khi thanh toán thành công sẽ được gắn một mã định danh duy nhất
(\texttt{ticketId}) và sinh kèm \textbf{QR code} trỏ đến trang thông tin
vé công khai (\texttt{/ticket-info/:ticketId}). Người dùng có thể xem, tải
ảnh QR về thiết bị hoặc nhận trực tiếp qua email. Ban tổ chức hoặc nhân
viên tại sự kiện có thể quét mã để kiểm tra tính hợp lệ của vé và tra cứu
đầy đủ thông tin người mua ở dạng HTML hoặc JSON thông qua API công khai
không cần xác thực.

\subsubsection*{(5) Hệ thống email tự động}

\indent \indent Nhóm xây dựng một email service hoàn chỉnh với nhiều loại thông báo: xác
nhận đặt vé (kèm QR từng vé inline trong email), cấp lại mật khẩu, và các
thông báo nghiệp vụ khác. Toàn bộ template email được thiết kế bằng HTML
thuần với logo thương hiệu nhúng dạng CID, tương thích với hầu hết các
ứng dụng email phổ biến. Dữ liệu động được escape HTML trước khi render
để phòng chống tấn công XSS qua email.

\subsubsection*{(6) Gợi ý sự kiện thông minh -- AI Chatbot}

\indent \indent Nhóm tích hợp \textbf{Google Gemini 2.5 Flash} làm trợ lý ảo tư vấn sự
kiện cho người dùng. Chatbot được cung cấp context động là danh sách tối
đa 30 sự kiện đang mở bán (bao gồm tên, thể loại, thời gian, địa điểm,
giá vé từng hạng), giúp trả lời các câu hỏi như ``có sự kiện nào rẻ dưới
500k không?'' hay ``sự kiện âm nhạc nào cuối tuần này?'' một cách tự nhiên
bằng tiếng Việt.

\subsubsection*{(7) Pipeline dữ liệu cho mô hình gợi ý ML}

\indent \indent Hệ thống xây dựng pipeline thu thập và xử lý dữ liệu tương tác người dùng
(click, purchase theo từng cặp user--event) thông qua model \texttt{Interaction}.
Dữ liệu được aggregate trong MongoDB (tính điểm tương tác, mã hóa đặc trưng
về tuổi, giới tính, thể loại sự kiện, vị trí địa lý, ngân sách trung bình)
và xuất ra file CSV, tự động upload lên \textbf{HuggingFace Hub} để phục vụ
việc huấn luyện mô hình gợi ý ngoại tuyến (\textit{offline training}).

\subsubsection*{(8) Dashboard thống kê}

\indent \indent Ban tổ chức có trang thống kê riêng với dữ liệu doanh thu, số vé bán được
theo sự kiện. Quản trị viên có dashboard tổng quan với các chỉ số: tổng số
sự kiện, tổng người dùng, tổng ban tổ chức, số liệu tăng trưởng trong tháng
(sự kiện mới, người dùng mới, ban tổ chức hoạt động). Dữ liệu thống kê
được lưu dạng snapshot định kỳ trong collection \texttt{statistics}, cho
phép theo dõi xu hướng theo thời gian.

\subsubsection*{(9) Đánh giá và phản hồi}

\indent \indent Hệ thống hiện thực chức năng đánh giá sự kiện với ràng buộc nghiệp vụ chặt
chẽ: chỉ người dùng đã mua vé và thanh toán thành công mới được đánh giá,
và chỉ sau khi sự kiện đã kết thúc. Mỗi người dùng chỉ được đánh giá một
lần cho mỗi sự kiện (ràng buộc unique index). Điểm đánh giá trung bình của
ban tổ chức được tự động cập nhật mỗi khi có đánh giá mới thông qua
MongoDB Aggregation Pipeline.

\subsubsection*{(10) Bảo mật tổng thể}

\indent \indent Toàn bộ các cơ chế bảo mật đã được hiện thực như mô tả trong Phần~1, bao
gồm: JWT với kiểm tra trạng thái tài khoản trong DB, hash mật khẩu bcrypt,
CORS whitelist, HTML escaping, DOMPurify, bảo vệ cron endpoint bằng secret
header, giới hạn số lượng bản ghi trả về và lọc input MongoDB để phòng
chống injection.

% -----------------------------------------------------------
\subsection{Hạn chế và hướng phát triển tương lai}

\indent \indent Mặc dù hệ thống đã đạt được những kết quả tích cực nêu trên, nhóm nhận thức
rõ rằng vẫn còn tồn tại một số hạn chế xuất phát từ điều kiện thời gian,
quy mô nhóm và phạm vi của đồ án tốt nghiệp. Các hạn chế này đồng thời mở
ra những hướng phát triển cụ thể cho giai đoạn tiếp theo.

\subsubsection*{Hạn chế hiện tại}

\begin{enumerate}[leftmargin=*, label=\textbf{\arabic*.}]

    \item \textbf{Mô hình gợi ý ML chưa được triển khai ở tầng phục vụ
    (inference serving).}
    Pipeline thu thập dữ liệu và upload lên HuggingFace đã được xây dựng
    đầy đủ, tuy nhiên việc tích hợp mô hình đã huấn luyện trở lại vào hệ
    thống để phục vụ gợi ý theo thời gian thực (real-time inference) cho
    từng người dùng chưa được hoàn thiện trong giai đoạn này. Hiện tại,
    chức năng gợi ý sự kiện chủ yếu dựa vào Gemini Chatbot với context
    tĩnh.

    \item \textbf{Chưa có bộ kiểm thử tự động (automated testing).}
    Cả backend lẫn frontend đều chưa có unit test hay integration test.
    Script \texttt{test} trong \texttt{package.json} của backend chỉ
    là placeholder (\texttt{echo "Error: no test specified"}). Điều này
    làm tăng rủi ro regression khi codebase phát triển thêm và khó đảm
    bảo độ bao phủ kiểm thử theo tiêu chuẩn công nghiệp.

    \item \textbf{Chưa có cơ chế rate limiting và throttling.}
    Hệ thống chưa tích hợp middleware giới hạn tần suất request (như
    \texttt{express-rate-limit}), khiến các endpoint nhạy cảm như đăng
    nhập, đăng ký, gửi email có thể bị lạm dụng bằng tấn công brute-force
    hoặc spam request với số lượng lớn.

    \item \textbf{Cơ chế reset mật khẩu chưa theo chuẩn token-based.}
    Luồng quên mật khẩu hiện tại sinh một mật khẩu ngẫu nhiên và gửi
    thẳng qua email. Cách tiếp cận này kém an toàn hơn so với chuẩn
    công nghiệp sử dụng token có thời hạn (OTP hoặc signed URL), vì
    mật khẩu mới bị truyền qua kênh email vốn không phải kênh bảo mật
    tuyệt đối.

    \item \textbf{BullMQ chưa được triển khai đầy đủ.}
    Mặc dù \texttt{bullmq} và \texttt{redis} đã được khai báo trong
    \texttt{package.json}, việc xử lý email và job nặng hiện vẫn chạy
    trực tiếp trong request cycle thay vì thông qua hàng đợi. Kiến trúc
    job queue chưa được kích hoạt hoàn toàn.

    \item \textbf{Chọn chỗ ngồi tương tác (interactive seat map) chưa
    được hiện thực hóa đầy đủ.}
    Model \texttt{Event} đã có trường \texttt{seatImgUrl} để lưu sơ đồ
    ghế và \texttt{TicketClass} hỗ trợ loại ghế \texttt{reserved}, nhưng
    giao diện chọn ghế tương tác trực quan (kéo-thả, tô màu theo trạng
    thái) chưa được phát triển. Người dùng hiện chưa thể tự chọn ghế cụ
    thể trên sơ đồ.

    \item \textbf{Chưa hỗ trợ cập nhật trạng thái vé theo thời gian thực.}
    Khi số lượng vé thay đổi (có người vừa mua hoặc cron job thu hồi vé),
    giao diện người dùng chỉ cập nhật khi họ tải lại trang, không có cơ
    chế đẩy dữ liệu theo thời gian thực (WebSocket / Server-Sent Events).
    Điều này có thể dẫn đến trải nghiệm kém trong các sự kiện có lượng
    người mua đồng thời cao.

    \item \textbf{Ràng buộc một tài khoản -- một ban tổ chức.}
    Hiện tại mỗi tài khoản \texttt{organizer} chỉ được liên kết với một
    thực thể Organizer duy nhất. Hệ thống chưa hỗ trợ mô hình tổ chức
    có nhiều thành viên (team-based organizer) hay phân quyền chi tiết
    trong nội bộ một ban tổ chức.

    \item \textbf{Chưa có ứng dụng di động.}
    Hệ thống hiện chỉ là ứng dụng web. Giao diện được tối ưu cho màn
    hình desktop và responsive cơ bản cho mobile browser, nhưng chưa có
    ứng dụng native (iOS/Android) để cung cấp trải nghiệm người dùng
    đầy đủ trên thiết bị di động, đặc biệt cho chức năng quét QR check-in.

    \item \textbf{Hệ thống thiếu kho mã giảm giá.}
    Hiện tại hệ thống chưa có cơ chế quản lý và mã giảm giá hợp lệ cho người dùng khi thanh toán vé do chưa thể làm việc được với nhà quản lý nền tảng thanh toán chính hiện tại là PayOs. 
    Việc này hạn chế khả năng triển khai các chiến dịch khuyến mãi và tăng cường trải nghiệm người dùng.

\end{enumerate}

\subsubsection*{Hướng phát triển tương lai}

\begin{enumerate}[leftmargin=*, label=\textbf{\arabic*.}]

    \item \textbf{Hoàn thiện và triển khai mô hình gợi ý cá nhân hóa.}
    Bước tiếp theo quan trọng nhất là hoàn thiện vòng lặp ML: huấn luyện
    mô hình collaborative filtering hoặc content-based filtering trên
    dataset đã thu thập tại HuggingFace, sau đó triển khai API inference
    (có thể dùng HuggingFace Inference Endpoints hoặc một microservice
    riêng) để tích hợp vào luồng gợi ý sự kiện theo người dùng cụ thể
    thay vì context chung như hiện tại.

    \item \textbf{Xây dựng bộ kiểm thử tự động và CI/CD pipeline.}
    Cần bổ sung unit test cho các controller và service quan trọng
    (thanh toán, đặt vé, xác thực) bằng Jest (backend) và React Testing
    Library (frontend), kết hợp tích hợp liên tục (CI) thông qua
    GitHub Actions để tự động kiểm thử mỗi khi có commit mới.

    \item \textbf{Tích hợp rate limiting và nâng cấp bảo mật.}
    Bổ sung \texttt{express-rate-limit} cho các endpoint nhạy cảm, nâng
    cấp luồng reset mật khẩu sang cơ chế token có thời hạn (6 chữ số
    OTP hoặc signed URL JWT), và bổ sung CAPTCHA cho form đăng ký/đăng
    nhập.

    \item \textbf{Kích hoạt BullMQ và tách worker process.}
    Triển khai đầy đủ kiến trúc job queue với BullMQ -- tách worker
    process gửi email và xử lý task nặng ra khỏi HTTP server process
    chính. Điều này giúp server phản hồi request nhanh hơn và tăng
    khả năng chịu tải khi lưu lượng tăng đột biến.

    \item \textbf{Phát triển chức năng chọn ghế tương tác.}
    Xây dựng component SVG/canvas cho phép người dùng trực quan chọn
    ghế theo sơ đồ, với màu sắc thể hiện trạng thái ghế (còn trống,
    đã đặt, đang giữ chỗ). Tích hợp WebSocket để cập nhật trạng thái
    ghế theo thời gian thực trong quá trình nhiều người cùng đặt.

    \item \textbf{Cập nhật thời gian thực với WebSocket.}
    Tích hợp Socket.IO hoặc Server-Sent Events để đẩy cập nhật số
    lượng vé còn lại, trạng thái sự kiện và thông báo hệ thống đến
    người dùng mà không cần tải lại trang, đặc biệt quan trọng trong
    các sự kiện có lưu lượng mua vé cao.

    \item \textbf{Phát triển ứng dụng di động.}
    Xây dựng ứng dụng React Native tái sử dụng logic từ các service
    TypeScript hiện có, tập trung vào các tính năng cốt lõi: duyệt và
    tìm kiếm sự kiện, đặt vé, hiển thị và quét QR check-in, nhận
    thông báo push khi sự kiện sắp diễn ra.

    \item \textbf{Nâng cấp hệ thống phân tích và báo cáo.}
    Bổ sung dashboard analytics chi tiết hơn cho ban tổ chức: biểu đồ
    doanh thu theo thời gian, phân tích nhân khẩu học người mua vé,
    tỷ lệ hoàn vé, và xuất báo cáo PDF/Excel. Tích hợp các chỉ số
    tương tác thực tế (tỷ lệ chuyển đổi click-to-purchase) vào
    dashboard để ban tổ chức tối ưu chiến lược bán vé.

    \item \textbf{Mở rộng tích hợp cổng thanh toán và phát triển kho mã giảm giá.}
    Bổ sung hỗ trợ các phương thức thanh toán phổ biến khác tại Việt
    Nam như VNPay, MoMo và ZaloPay, đồng thời xây dựng luồng hoàn tiền
    (refund) tự động khi sự kiện bị hủy.\\
    Phát triển hệ thống quản lý mã giảm giá (coupon) với các loại chiết khấu: 
    theo phần trăm, theo số tiền cố định, hoặc miễn phí vận chuyển (nếu có). Mã giảm giá sẽ có điều kiện áp dụng (tối thiểu hóa đơn, chỉ áp dụng cho sự kiện cụ thể, giới hạn số lượng sử dụng) 
    và được tích hợp vào luồng thanh toán.

    \item \textbf{Hỗ trợ đa ngôn ngữ.}
    Mở rộng hệ thống hỗ trợ tiếng Anh ngoài tiếng Việt, sử dụng thư
    viện \texttt{i18next} hoặc \texttt{react-intl}, nhằm phục vụ cả
    người dùng quốc tế và ban tổ chức sự kiện có yếu tố nước ngoài.

\end{enumerate}

\subsubsection*{Kết luận}

\indent \indent Nhìn lại toàn bộ quá trình thực hiện, hiện nhóm
đã thành công trong việc xây dựng một hệ thống đặt vé sự kiện trực tuyến
đầy đủ chức năng, tích hợp các công nghệ hiện đại bao gồm xác thực JWT,
thanh toán điện tử, vé QR điện tử, email tự động và trí tuệ nhân tạo,
được triển khai và vận hành thực tế trên nền tảng cloud. Các hạn chế hiện
tại chủ yếu là những tính năng nâng cao mà nhóm đã có kế hoạch kiến trúc
rõ ràng nhưng chưa đủ thời gian hoàn thiện trong phạm vi đồ án. Những
hướng phát triển được đề xuất đều có cơ sở từ kiến trúc hệ thống đã xây
dựng, tạo nền tảng vững chắc để nhóm hoặc các thành viên tiếp theo có thể
phát triển liên tục sau khi kết thúc đồ án.